\section{The dyad kernel and complete shape parameterization}\label{sec:dyadfull}

The periodic weights in D can be generated without solving an inverse
homometry problem. We give the kernel theorem and the recovery argument
explicitly. Throughout this section, functions have integer coefficients
on $C_n$; translation acts by $x^aW(t)=W(t-a)$.

\begin{lemma}[Nonnegative periodic decomposition]\label{lem:periodic}
Let $W\colon C_n\to\Z_{\ge0}$ and $a,b\in C_n$. Then
\[
 (1-x^a)(1-x^b)W=0
 \quad\Longleftrightarrow\quad
 W=U+V,
\]
where $U,V$ are nonnegative integer functions satisfying
$x^aU=U$ and $x^bV=V$. If $W$ has total mass $m$, then
\[
 m=\ord(a)\sum_{C_n/\langle a\rangle}U
   +\ord(b)\sum_{C_n/\langle b\rangle}V,
\]
where the sums take one value per coset.
\end{lemma}
\begin{proof}
The reverse direction follows by annihilating the two summands. For the
forward direction use the characters $\chi_k(t)=\exp(2\pi i kt/n)$,
$0\le k<n$. Translation is diagonal on these characters, and a product
eigenvalue $(1-\chi_k(-a))(1-\chi_k(-b))$ vanishes precisely when at least
one of its factors does. Hence the complex kernel of the product is the
sum of the two periodic subspaces. Taking real parts gives real periodic
functions $U_0,V_0$ with $W=U_0+V_0$.

This alone does not give nonnegative integral summands. Put
$H=\langle a\rangle$, $J=\langle b\rangle$ and $L=H+J$, and work
in one coset of $L$. List its $H$-cosets by $i$ and its $J$-cosets by $j$.
Every $H$-coset meets every $J$-coset: the difference between their
representatives belongs to $H+J$. All intersections are nonempty, have
the same size $|H\cap J|$, and partition the component. The function
$W=U_0+V_0$ is constant on an intersection; call its integer value $w_{ij}$.
It satisfies the rectangle identities
\[
 w_{ij}-w_{ij_0}=w_{i_0j}-w_{i_0j_0}
 \qquad\text{for all }i,j.
\]
Choose $j_0$ and set $m_0=\min_i w_{ij_0}$,
\[
 u_i=w_{ij_0}-m_0,\qquad
 v_j=w_{i_0j}-w_{i_0j_0}+m_0.
\]
These numbers are integral, $u_i\ge0$, and the rectangle identities give
$v_j=\min_i w_{ij}\ge0$ and $w_{ij}=u_i+v_j$. Assign $u_i$ to its whole
$H$-coset and $v_j$ to its whole $J$-coset. Repeat on each $L$-coset.
The resulting $U,V$ are the required nonnegative integral functions.
Summing their constant values along cosets proves the mass formula.
\end{proof}

\begin{proposition}[Exactness of D within the dyad shape]\label{prop:dyadshape}
Suppose independent rigid motions align a six-pair to
\[
 X=C\sqcup\{0,a+b\},\qquad Y=C\sqcup\{a,b\},
\]
where $|C|=4$ and $R=\{0,a,b,a+b\}$ consists of four distinct elements
outside $C$. The pair is homometric if and only if it has the periodic
weights and orbit selections specified in construction D.
\end{proposition}
\begin{proof}
Write $\rho(C)=x^{a+b}C^*$ and $K=C+\rho(C)+1_R$.
The expansion \eqref{eq:dyad} and invertibility of $x^{-a-b}$ show that
homometry is equivalent to $(1-x^a)(1-x^b)K=0$. The function $K$ is
nonnegative, integral, has mass twelve, and is fixed by $\rho$.
Its coefficients are one on $R$ and belong to $\{0,1,2\}$ off $R$.
At a $\rho$-fixed point off $R$ its coefficient is twice the coefficient
of $C$, and therefore belongs to $\{0,2\}$. Lemma~\ref{lem:periodic}
supplies exactly the required weights.

Conversely, admissible weights give a $K$ with these properties. On each
two-point $\rho$-orbit, choose no point, one point, or both points according
as its coefficient is zero, one, or two. At a fixed point choose it
precisely when its coefficient is two. This produces a binary set $C$
off $R$ with $C+\rho(C)=K-1_R$. Summing coefficients gives
$2|C|=12-4=8$, so its cardinality is four. The dyad identity now gives
homometry. The same selection rule recovers every $C$ in the forward
direction. The final $\TI$ test decides whether it is a Z-pair.
\end{proof}

If $\ord(a)>12$ then its periodic summand must vanish, because a nonzero
nonnegative integer weight would alone contribute more than twelve.
Thus $K$ is $b$-periodic and $\ord(b)$ divides twelve. If both orders
exceed twelve, the construction is impossible. These are consequences
of the exact mass formula, not search caps. This proposition classifies
the displayed exchange shape; it does not assert that every six-pair
can be aligned to that shape.

\subsection{Coefficient proof of half-per-coset complementation}
Let $H$ have even order $h$, and let $A$ contain $h/2$ points in each
of $m=12/h$ selected $H$-cosets. Write $S$ for their full union. For
any $t\in C_n$, let $N_t$ count the ordered pairs of selected cosets
$(K,J)$ with $K=t+J$. Then direct coefficient counting gives
\[
 (SS^*)(t)=hN_t,\qquad
 (AS^*)(t)=(SA^*)(t)=\tfrac h2N_t.
\]
For example, for each such pair $(K,J)$ exactly $h/2$ choices of a point
of $A\cap K$ have their difference-$t$ partner in $J$. Consequently
$AS^*+SA^*=SS^*$, and expansion of $(S-A)(S-A)^*$ proves L7.
The equality of the two cross coefficients here follows from the
half-coset hypothesis, rather than from a general symmetry assumption
on $A$. The number of selected cosets is integral precisely for
$h\in\{2,4,6,12\}$ in the six-point construction.

\section{The exact weighted real-line obligation}\label{sec:linefull}

The following exposition states the complete mathematical formula whose
unsatisfiability is used in BF. It is included so that the real-domain
claim can be evaluated independently of a program's normal-form names.
A six-atom multiset is a list of six real positions with repetitions
allowed; the fifteen unordered pairs refer to distinct labelled atoms,
even when their positions coincide.

\begin{lemma}[Diameter normalization]\label{lem:linenorm}
Every noncongruent homometric pair of six-atom real multisets has positive
common diameter $L$. After independent translations, reflections and
division by $L$, followed if necessary by interchange of endpoints, it
has the form
\[
 A=(0,a_1,a_2,a_3,r,1),\qquad B=(0,b_1,b_2,b_3,r,1),
\]
with weakly increasing coordinates, $r\ge1-a_1$, $r\ge1-b_1$, and
$(a_1,a_2,a_3)<_{\rm lex}(b_1,b_2,b_3)$. The two lists are not related
by reflection about $1/2$.
\end{lemma}
\begin{proof}
The largest labelled pair distance is the diameter, so homometry gives
the same diameter for both lists. Diameter zero would give two congruent
lists with all six atoms at one point. Anchor the minima at zero and
the maxima at one. Remove one labelled endpoint-edge occurrence of
length one from the common multiset of distances. In either list the
largest remaining value is $\max(a_4,1-a_1)$: both those endpoint
distances remain, and every other edge is no larger than one of them.
This argument still holds if an endpoint is repeated or the remaining
maximum is also one. Reflect each list independently to put its shared
remaining maximum at coordinate four. The two lists are unequal, since
equality would be a congruence; lexicographic interchange is therefore
available. Reflection is excluded by the original noncongruence.
\end{proof}

Put $z=(a_1,a_2,a_3,r,1,b_1,b_2,b_3,r,1)$. The counterexample formula
$\Phi(z)$ is the conjunction of the following predicates:
\begin{enumerate}
\item Each displayed list is weakly increasing and the two inequalities
$r\ge1-a_1$, $r\ge1-b_1$ hold.
\item The strict lexicographic inequality of Lemma~\ref{lem:linenorm}
holds, and
$\neg\bigwedge_{i=1}^{4}(A_i=1-B_{5-i})$ holds.
\item With $d_{ij}=A_j-A_i$ and $e_{kl}=B_l-B_k$, for every labelled
$i<j$,
\[
 \sum_{k<l}[d_{kl}=d_{ij}]
   =\sum_{k<l}[e_{kl}=d_{ij}].
\]
Here $[P]$ is one if $P$ holds and zero otherwise.
\item The vector $z$ belongs to neither of the two linear planes given
by the two height lists displayed in BF.
\end{enumerate}
The indices in the last condition run through the ten coordinates after
omitting the two anchored zeros. If a template has coefficient pairs
$(u_k,v_k)$, choose $i,j$ with $\Delta=u_iv_j-v_iu_j\ne0$. Membership
in its plane is the conjunction, for all ten indices $k$, of
\[
 \Delta z_k=(u_kv_j-v_ku_j)z_i+(-u_kv_i+v_ku_i)z_j.
\]
Thus nonmembership is the negation of a specified finite conjunction;
there is no bounded parameter grid or rationality condition.

\begin{lemma}[Encoding equivalence]\label{lem:lineencode}
$\Phi$ is satisfiable over the reals if and only if there is a normalized
noncongruent homometric six-atom pair outside both displayed planes.
\end{lemma}
\begin{proof}
Only the distance-count and plane predicates need explanation. The
distance equations give equality of multiplicity for every value occurring
among the fifteen $A$ edges, including zero. The required multiplicities
sum to fifteen, so there is no room for a further $B$ value. Conversely,
equal multisets satisfy every equation. Equalities among labelled atoms
are retained; the directed autocorrelation adds six diagonal zeros and
twice the unordered zero distances to both sides. Cramer's rule shows
that the displayed plane equations are exactly the existence of real
parameters $p,q$ with $z_k=u_kp+v_kq$. The other predicates are the
normalization and exclusion conditions already proved.
\end{proof}

\paragraph{Finite solver obligation U.}
The exact formula $\Phi$ is unsatisfiable. This is the saved exact
linear-real-arithmetic obligation checked by the weighted builder and
a separate consecutive-gap/occurrence-bijection encoding. Its truth at
all real positions is a solver dependency of this computer-assisted
proof. A bounded integer recount does not establish U. Saved Z3 and
cvc5 results, input hashes and internal proof-checking settings are
specified in the evidence inventory; their proof objects have not been
translated into an independent general-purpose proof assistant.

Accepting U, Lemmas~\ref{lem:linenorm} and~\ref{lem:lineencode} put every
pair in one of the two planes. Recovering the parameters in the first
plane gives $p=A_5-A_4$, $q=A_2$. The weak consecutive gaps force
$p\ge0$ and $q\ge3p$; $p=0$ makes the lists equal. In the second plane,
$p=A_4-A_3$, $q=A_1+p$; the gaps force $p>0$,
$3p/2\le q\le2p$, and the upper endpoint makes the lists equal.
Thus the exact domains are the relaxed chambers stated in BF. Their
formal positive-distance identities prove soundness even on the allowed
repeated-coordinate boundaries. For completeness their consecutive gap
lists are, respectively,
\begin{align*}
 A_1 &: (q-2p,2p,q-3p,q+p,p),\\
 B_1 &: (q-p,2p,q,q-3p,p),\\
 A_2 &: (q-p,2p,2q-3p,p,2p-q),\\
 B_2 &: (p,2q-2p,3p-q,2q-3p,2p-q).
\end{align*}
The first coordinates after zero in the first pair differ by $p>0$.
Reflecting $A_1$ makes that coordinate $p$, whereas $B_1$ has
$q-p\ge2p$. In the second pair the two coordinates are $q-p$ and $p$,
which differ because $q<2p$; the reflected $A_2$ has coordinate
$2p-q\le p/2<p$. Anchored equal-diameter lists can be congruent only by
equality or this reflection. Thus both forms are noncongruent, including
their allowed repeated-coordinate boundaries. Integer input coordinates recover integer
$p,q$ after undoing the diameter normalization: the formulas use only
integer point differences. For distinct six-sets the lower bounds become
strict. This is the full logical passage from U to Iw and I.

\section{Integer shadows and the large-prime boundary}\label{sec:shadowfull}

\begin{proposition}\label{prop:shadowexact}
Assuming the line obligation U, a cyclic Z-pair is an integer shadow
if and only if representatives are an admissible modular Bloom pair.
\end{proposition}
\begin{proof}
Apply the distinct-point line theorem I to any integer lifts and reduce
its integer rigid motions and recovered integer parameters modulo $n$.
For the converse lift the two modular parameters to any integers $p,q$.
Six distinct residues in each formula force six distinct integer terms.
The formal Laurent identity for B is valid in $\Z[x,x^{-1}]$, so those
integer sets are homometric. An integer congruence would reduce to a
cyclic $\TI$ congruence; the latter has been excluded. This also explains
why exhaustive testing of the $n^2$ modular parameter pairs, with rigid
alignments, is an exact shadow test despite unbounded integer lifts.
\end{proof}

\begin{corollary}\label{cor:largeprime}
If every prime divisor of $n$ is greater than 131, every cyclic six-point
Z-pair is an admissible modular Bloom image, under the same line-solver
qualification.
\end{corollary}
\begin{proof}
Take its universal presentation $G_M=\Z^d\oplus T$ and map
$f\colon G_M\to C_n$. The torsion image has order dividing both $n$
and $|T|\le135$. The smallest prime above 131 is 137, so this image
is trivial. The pair factors through the free quotient. Its six vertices
in each list remain distinct, since their images in $C_n$ are distinct.
Here $d>0$, because a finite presentation with trivial image cannot
contain a six-set.

Choose a basis of the free quotient and any integer representatives
$c_1,\ldots,c_d$ of its images under $f$. The integer functional
$\lambda(w)=\sum_j c_jw_j$ reduces modulo $n$ to $f$. Any within-endpoint
difference $w$ has nonzero image under $f$, so $\lambda(w)$ is a nonzero
integer. Thus $\lambda$ is injective on each six-element endpoint, with
no genericity assumption. Homomorphic transport gives integer homometric
six-sets reducing to the original cyclic pair.
Proposition~\ref{prop:shadowexact} finishes.
\end{proof}

This sufficient boundary is not an enumeration formula. In particular,
deducing $(p-1)(p-11)/12$ at primes requires a separate proof of parameter
admissibility, all identifications between parameter images, and all
exceptional stabilizers. That counting problem is not settled merely by
the large-prime corollary or by matching finite numerical values.

\subsection{What the rigid-template certificates additionally prove}
The soundness of an R edge requires only equality of its fifteen unsigned
differences modulo its listed $q$. The stronger claim that a faithful seed
image is purely cyclic uses an additional, pattern-restricted statement.
We give its exact finite contract to keep these two uses separate.

For each row $(q,X,Y)$ of Table~\ref{tab:rigid}, order its six coordinates
as printed and exhaust all bijections of $\mathcal E$ and all signs
compatible with those coordinates modulo $q$. Compatibility means the
row relation~\eqref{eq:formal-matching-row} evaluates to zero modulo $q$;
both signs are retained for an involution. The saved arithmetic witnesses
must, for every such matching and with no omitted or duplicate label list,
supply integer unimodular $U,V$ satisfying
\[
 UMV=\operatorname{diag}(1,1,1,1,1,1,1,1,1,q),
\]
with the five additional rows zero. If $w$ is the last column of $V$ and
$v_{\rm seed}$ is the ten-coordinate anchored seed vector, also require
\[
 \gcd(w_1,q)=1,\qquad w_i\equiv w_1(v_{\rm seed})_i\pmod q
 \quad(1\le i\le10).
\]
The seed's first nonanchored coordinate is one. These predicates are
checked by integer multiplication, determinants and congruences. The
numbers of compatible matchings, in table order, are
$192,96,48,48,48,48,16,12,8,8,4,4,4$, totaling $536$.

\begin{proposition}[Rigidity restricted to the seed matching patterns]
Under the preceding finite verification, every realization of a compatible
seed matching in any abelian group is the displayed seed multiplied by
one element $t$ with $qt=0$. Every faithful cyclic image of each listed
seed is purely cyclic.
\end{proposition}
\begin{proof}
The unimodular diagonal identity identifies its universal presentation
with $C_q$. All old vertices are multiples of the generator, with
coefficients $w_i$. Replacing the generator by $w_1$ times itself is
invertible modulo $q$ and gives exactly $v_{\rm seed}$. Thus every
realization has the asserted form. The allowed orders in the table are
checked by reducing both seed lists to each divisor of $q$, requiring
six distinct coordinates and different $\TI$ classes. By the subgroup
equivalence argument in Lemma~\ref{lem:inflate}, that check is valid in
every ambient cyclic group.

If the seed itself had homometric six-point integer lifts, their equal
absolute edge lengths would give one of the exhausted signed matchings
after reduction modulo $q$. Its anchored integer coordinates would lie
in the kernel of a rank-ten $M$, hence would all be zero, contradicting
six distinct points. Now let a faithful image lie in $C_n$, with
$t=(n/q)u$ and $\gcd(u,q)=1$. Anchor hypothetical integer lifts at their
points mapping to zero. All their coordinates are divisible by $n/q$.
Divide by $n/q$, then multiply by any integer inverse of $u$ modulo $q$.
These operations preserve integer homometry and distinctness and would
give integer lifts of the original seed, a contradiction.
\end{proof}

The proposition concerns the listed signed matching patterns, not all
six-point pairs. Proper images can admit additional matching patterns;
their pure cyclicity does not follow from the faithful-image argument.
In particular, this verifies the 21-EDO benchmark's pure cyclicity without
assuming that a cyclic factor flip is an integer-line factorization.
